\scanpage{88}

\deleted{$G'=G_{S'}$ soit diagonalisable, i.e. de la forme $D_{S'}(M)$. Il est immédiat
qu'on peut se ramener au cas des $G'$, donc on peut supposer $G$ diagonalisable, donc il provient du groupe absolu $D_{\underline{Z}}(M)$ sur $\mathrm{Spec}(\underline{Z})$ par changement de base $S\quad\mathrm{Spec}(\underline{Z})$. Il suffit maintenant de conjuguer a) et 4.4.}

\deleted{\uline{Corollaire} 4.\handwritten{\unsure{8|7|9}}. \quad a) Soient $u,v:H\quad G$ deux homomorphismes de $S$-préschémas en groupes, avec $H$ diagonalisable de type fini, et \unsure{supposons|supposon} que pour tout entier $n\quad 0$, les restrictions de $u$ et $v$ à $ {}_nH$ sont identiques,. Alors $u\unsure{=}v$. \quad b) Soient $H_1,H_2$ deux sous-groupes de type multiplicatif et de type fini d'un préschéma en groupes $G$, et supposons que pour tout entier $n\quad 0$, on ait $ {}_nH_1=\deleted{\unsure{n}}{}_nH_2$, alors $H_1=H_2$.}

\deleted{La première assertion résulte de la seconde, par considération des
sous-groupes $H_1,H_2$ de $H\mathbin{\mathrm{x}}G$ graphes de $u$ et $v$. Pour prouver b), soit
$H=H_1\quad H_2$, c'est un sous-préschéma \inserted{en group\unsure{e|é}s} de \deleted{$H_1\unsure{H|M}$} $H_i$ $(i=1,2)$, il faut montrer qu'il est identique à $H_i$, or l'hypothèse signifie qu'il majore les $ {}_nH_i$. On est donc ramené à prouver le}

\deleted{\uline{Corollaire} 4.\handwritten{9}. \quad Soient $G$ un $S$-groupe de type multiplicatif et de type fini, et $H$ un sous-préschéma en groupes qui majore les $ {}_nG$, $n\quad 0$. Alors $H=G$.}

\deleted{En vertu de 4.6. on est ramené à prouver que le sous-préschéma
$\unsure{\mathbb H|H}$ est \uline{fermé}, ou encore qu'il est \deleted{\unsure{ident}} égal ensemblistement à $G$. Cela nous ramène au cas où $S$ est le spectre d'un corps, mais alors tout sous-préschéma en groupes de $G$ est fermé (Exp V), d'où la conclusion.}

\deleted{\uline{Remarque} 4.9. \quad Sous les conditions de 4.6., soit $m$ un entier $\quad 0$ qui a les propriétés suivantes: pour tout $s\quad S$, $m$ n'est pas \deleted{\unsure{la}} \inserted{une} puissance de la caractéristique de $k(s)$, et si $G_s$ est de type $M$, les diviseurs premiers de la torsion de $M$ divisent $m$. (NB cette deuxième condition}

% archive sheet number: 87
